\documentclass[11pt,a4paper]{ctexart}
\usepackage{amsmath,amssymb,amsthm}
\usepackage{mathtools}
\usepackage{geometry}
\geometry{margin=2.5cm}
\usepackage{hyperref}
\usepackage{enumitem}
\usepackage{cleveref}

\newtheorem{theorem}{定理}[section]
\newtheorem{lemma}[theorem]{引理}
\newtheorem{corollary}[theorem]{推论}
\newtheorem{proposition}[theorem]{命题}
\newtheorem{definition}[theorem]{定义}
\newtheorem{remark}[theorem]{注记}
\newtheorem{openproblem}[theorem]{开放问题}
\newtheorem{conjecture}[theorem]{猜想}

\newcommand{\F}{\mathbb{F}}
\newcommand{\Z}{\mathbb{Z}}
\renewcommand{\d}{\mathsf{d}}
\DeclareMathOperator{\rank}{rank}
\DeclareMathOperator{\width}{width}
\DeclareMathOperator{\cutrank}{cutrank}
\DeclareMathOperator{\im}{im}
\DeclareMathOperator{\spanop}{span}
\DeclareMathOperator{\Gram}{Gram}
\DeclareMathOperator{\Supp}{Supp}
\DeclareMathOperator{\codim}{codim}
\DeclareMathOperator{\gr}{gr}
\DeclareMathOperator{\coker}{coker}
\newcommand{\obs}{\mathrm{obs}}

\begin{document}

\title{线性元胞自动机的有限时空可观测性：\\
关联分级、leading term 与代数障碍}
\author{}
\date{}
\maketitle

\begin{abstract}
研究有限域 \(\F_q\) (\(q=p^k\), \(p\) 素数) 上一维线性元胞自动机在有限时空窗口内的可观测性。
初始构型具有有限支撑，观测限制在与传播锥同步扩展的窗口内。
观测矩阵 \(\Theta_T\) 汇总 \(T\) 步内全部时空观测，其核维数 \(H_T=\dim\ker\Theta_T\) 称为观测缺陷。
本工作建立关联分级与 leading term 的统一框架，主要贡献如下。
\begin{enumerate}[label=(\arabic*)]
    \item 引入几何接触时间 \(\tau(i)\) 与几何过滤 \(G_t\)，证明几何核 \(G_T\) 维数由极值步幅 \(d_-,d_+\) 显式给出，构成缺陷的严格下界（命题~\ref{prop:geo_kernel}）。
    \item 构造核过滤 \(\{K_t\}\) 及其关联分级 \(\gr(V_0)=\bigoplus E_t\)。观测矩阵诱导分次映射 \(\gr(\Theta_T)=\bigoplus\theta_t\)，其中 \(\theta_t:E_t\to\F_q^{[-cT,cT]}\) 捕获首次非零观测（定理~\ref{thm:associated_graded}）。
    \item 定义形式首次观测的 leading term。证明：若各基向量的 leading term 两两不同，则 \(\gr(\Theta_T)\) 单，从而 \(\Theta_T\) 列满秩（定理~\ref{thm:leading_term_injective}）。该定理将极值满秩子矩阵与规则90等特例统一为同一结构。
    \item 引入碰撞图与 leading term 矩阵，证明一般情形的缺陷上界可由碰撞数给出，并将代数核完全归约为 leading term 碰撞所诱导的线性关系（命题~\ref{prop:collision_bound}）。
    \item 作为直接推论，严格证明当 \(r=1\) 且端部系数非零时 \(H_T=0\)（定理~\ref{thm:r1_full_obs}）。对 \(r\ge2\)，将未决问题明确拆分为几何版（无 Frobenius 共振时是否总有 \(H_T=0\)）与有限域版（Frobenius 稀疏化是否生成反例），并指出 leading term 碰撞分析是解决这些问题的核心路径。
\end{enumerate}
所有参数均显式可计算。已证部分逻辑闭合，未决部分以猜想与开放问题的形式明确标出。
\end{abstract}

\section{引言与问题设定}
设状态空间为有限域 \(\F_q\)（\(q=p^k\)，\(p\) 素数）。考虑一维线性规则
\begin{equation}\label{eq:local_rule}
    x_i^{(t+1)} = \sum_{j=-r}^{r} a_j x_{i+j}^{(t)},\qquad a_j\in\F_q,\; r\ge 1,
\end{equation}
全局演化由卷积算子 \(F:\F_q^{\Z}\to\F_q^{\Z}\) 给出：\((F x)_i = \sum_{j=-r}^r a_j x_{i+j}\)。

本文研究如下问题：已知初始构型具有有限支撑，且仅能通过时空窗口
\begin{equation}\label{eq:window}
    \Lambda_T = \{(t,i)\mid 0\le t\le T,\ |i|\le cT\},\qquad c\ge r,
\end{equation}
内的状态值进行观测，能否唯一恢复该初始构型？窗口随 \(T\) 线性扩展，恰与最大传播速度 \(r\) 相匹配，保证 \(\Lambda_T\) 内的状态完全由初始时刻落在
\begin{equation}
    A_T = \{i\in\Z : |i|\le (c+r)T\}
\end{equation}
上的构型决定。将初始构型限制在 \(V_0 = \{x\in\F_q^{\Z}:\Supp(x)\subseteq A_T\}\) 上，标准基为 \(\{e_i\}_{i\in A_T}\)。
定义观测投影 \(\pi:\F_q^{\Z}\to \F_q^{[-cT,cT]}\) 为限制映射，并令观测空间 \(W_T = \bigoplus_{t=0}^T \F_q^{[-cT,cT]}\)。
观测矩阵
\begin{equation}\label{eq:Theta_def}
    \Theta_T : V_0 \longrightarrow W_T,\qquad
    \Theta_T x = (\pi x,\ \pi F x,\ \dots,\ \pi F^T x),
\end{equation}
完全编码了有限时空窗口内的全部可用信息。
记 \(n_T = |A_T| = 2(c+r)T+1\)，\(R_T = \rank \Theta_T\)，\(H_T = \dim\ker\Theta_T = n_T - R_T\)，称为观测缺陷。

上述设置可等价地用 Laurent 多项式表述。将 \(x\) 视为 Laurent 多项式 \(X(z)=\sum_i x_i z^i\)，则
\[
F^t x \;\longleftrightarrow\; P(z)^t X(z),\qquad P(z)=\sum_{j=-r}^{r} a_j z^j.
\]
观测条件 \(\pi F^t x = 0\) 等价于
\begin{equation}\label{eq:moment}
    [z^j]\, P(z)^t X(z) = 0,\qquad \forall\,|j|\le cT,\; 0\le t\le T.
\end{equation}
因此 \(\ker\Theta_T\) 恰为满足上述矩方程的非零有限支撑 Laurent 多项式 \(X(z)\) 的集合。

\section{可观测性等价判据}
\begin{theorem}[可观测性等价判据]\label{thm:equivalent}
以下陈述等价：
\begin{enumerate}[label=(\roman*)]
    \item \(H_T=0\)（完全可观测）；
    \item \(\Theta_T\) 列满秩；
    \item \(\ker\Theta_T = \{0\}\)；
    \item 分离性：对任意非零 \(x\in V_0\)，存在 \(t\in\{0,\dots,T\}\) 使 \(\pi F^t x \neq 0\)。
\end{enumerate}
\end{theorem}
\begin{proof}
由秩–零化度定理直接可得。
\end{proof}

\section{几何障碍与几何核}
规则多项式的极值非零位移定义为
\begin{align*}
    d_- &= \max\{|j|: j<0,\ a_j\neq 0\}, \\
    d_+ &= \max\{j: j>0,\ a_j\neq 0\},
\end{align*}
若某方向无任何非零系数，则对应值取 \(0\)。显见 \(0\le d_-,d_+\le r\)。

\begin{lemma}[有限传播]\label{lem:finite_propagation}
设 \(x\in V_0\)。对任意 \(t\ge 0\)，
\[
\Supp(F^t x) \subseteq \Supp(x) + [-d_- t,\ d_+ t].
\]
特别地，对基向量 \(e_i\)，\(\Supp(F^t e_i)\subseteq [i-d_- t,\ i+d_+ t]\)。
\end{lemma}

\begin{definition}[几何接触时间]\label{def:tau}
对每个基索引 \(i\in A_T\)，定义其\emph{几何接触时间}
\[
\tau(i) = \min\{ t\in\{0,\dots,T\} : [i-d_-t,\; i+d_+t] \cap [-cT,cT] \neq \emptyset \}.
\]
若对所有 \(t\le T\) 交集为空，则记 \(\tau(i)=\infty\)。
显式公式为
\[
\tau(i)=
\begin{cases}
0, & |i|\le cT,\\[4pt]
\displaystyle\Big\lceil\frac{-cT-i}{d_+}\Big\rceil, & i< -cT \text{ 且 } d_+>0,\\[10pt]
\displaystyle\Big\lceil\frac{i-cT}{d_-}\Big\rceil, & i>cT \text{ 且 } d_->0,\\[10pt]
\infty, & \text{否则。}
\end{cases}
\]
若计算值超过 \(T\) 或对应方向 \(d=0\)，则 \(\tau(i)=\infty\)。
\end{definition}

\textbf{注：} \(\tau(i)\) 仅描述几何上首次可能落入窗口的时间，不保证该时刻观测值非零；由于代数抵消，实际首次非零观测时间可能大于 \(\tau(i)\) 或根本不存在。

记 \(I_\ell = \{i\in A_T : \tau(i)=\ell\}\)，\(\ell=0,\dots,T\)，以及 \(I_\infty = \{i:\tau(i)=\infty\}\)。得到直和分解
\[
V_0 = \Big( \bigoplus_{\ell=0}^{T} V_\ell \Big) \oplus V_\infty,\qquad
V_\ell = \spanop\{e_i : i\in I_\ell\},\; V_\infty = \spanop\{e_i : i\in I_\infty\}.
\]
\(\dim V_\ell = |I_\ell|\) 可显式计算（附录~\ref{sec:card_I}）。

\begin{definition}[几何过滤]\label{def:geo_filtration}
对 \(t=-1,0,\dots,T\) 定义
\[
G_t = \spanop\{ e_i : \tau(i) > t \},
\]
并约定 \(G_{-1}=V_0\)。显然 \(G_T = V_\infty\)。
\end{definition}

\begin{proposition}[几何过滤与核过滤的包含关系]\label{prop:geo_inclusion}
设核过滤 \(K_t = \bigcap_{s=0}^{t}\ker(\pi F^s)\)，\(K_{-1}=V_0\)。则对任意 \(t\ge -1\) 有
\[
G_t \subseteq K_t .
\]
特别地，\(G_T\subseteq \ker\Theta_T\)，且 \(\dim G_T = |I_\infty|\)。
\end{proposition}
\begin{proof}
若 \(e_i\in G_t\)，则 \(\tau(i)>t\)，故对所有 \(s\le t\)，区间 \([i-d_-s,i+d_+s]\) 与 \([-cT,cT]\) 不交，从而 \(\pi F^s e_i =0\)，即 \(e_i\in K_t\)。由线性张成即得 \(G_t\subseteq K_t\)。
\end{proof}

称 \(G_T\) 为\emph{几何核}，它完全由有限传播速度决定，不涉及任何代数抵消。商空间
\[
K_T^{\mathrm{alg}} = \ker\Theta_T / G_T
\]
称为\emph{代数核}。于是有正合列
\[
0 \longrightarrow G_T \longrightarrow \ker\Theta_T \longrightarrow K_T^{\mathrm{alg}} \longrightarrow 0.
\]

\begin{proposition}[几何核维数——缺陷下界]\label{prop:geo_kernel}
\[
\dim G_T = T(r-d_-) + T(r-d_+).
\]
若 \(a_{-r}=0\) 或 \(a_r=0\)，则 \(\dim G_T>0\)，从而 \(H_T>0\)。
\end{proposition}
\begin{proof}
\(I_\infty\) 由满足 \(i > cT + d_-T\) 或 \(i < -cT - d_+T\) 的索引构成。在 \(A_T = [-(c+r)T, (c+r)T]\) 内计数即得维数公式。
\end{proof}

\section{核过滤与关联分级}
尽管 \(\tau(i)\) 不能直接给出首次非零观测时间，但核过滤 \(K_t\) 精确刻画了可观测性的逐层结构。

\begin{definition}[首次可见时间]
对于非零 \(v\in V_0\)，定义其\emph{首次可见时间}
\[
\nu(v) = \min\{ t\in\{0,\dots,T\} : v\notin K_t \},
\]
若 \(v\in K_T\) 则记 \(\nu(v)=\infty\)。
\end{definition}

令 \(E_t = K_{t-1}/K_t\)（\(t=0,\dots,T\)），则 \(E_t\) 中的元素恰为在时刻 \(t\) 首次产生非零观测的自由度。由线性代数基本定理得向量空间同构
\begin{equation}\label{eq:decompose}
    V_0/K_T \cong \bigoplus_{t=0}^T E_t.
\end{equation}
\(\{K_t\}\) 是 \(V_0\) 上关于观测算子的标准降滤过，而 \(\gr(V_0)=\bigoplus_{t=0}^T E_t\) 为其关联分级对象。

观测矩阵 \(\Theta_T\) 自然地诱导出分级映射。对每个 \(t\)，定义
\[
\theta_t : E_t \longrightarrow \F_q^{[-cT,cT]},\qquad
\theta_t(v + K_t) = \pi F^t v.
\]
由于 \(v\in K_{t-1}\)，对于 \(s<t\) 有 \(\pi F^s v = 0\)，故 \(\theta_t\) 良定义且线性。分次映射
\[
\gr(\Theta_T) = \bigoplus_{t=0}^T \theta_t : \bigoplus_{t=0}^T E_t \longrightarrow \bigoplus_{t=0}^T \F_q^{[-cT,cT]} = W_T
\]
满足 \(\gr(\Theta_T)\circ q = \Theta_T\)，其中 \(q: V_0 \to \gr(V_0)\) 是商映射。这是滤过线性映射的标准关联分级构造。

\begin{theorem}[关联分级与观测矩阵]\label{thm:associated_graded}
以下陈述等价：
\begin{enumerate}[label=(\roman*)]
    \item \(\Theta_T\) 在 \(V_0/G_T\) 上单射（即代数核为零）；
    \item \(\gr(\Theta_T)\) 单射；
    \item 对每个 \(t=0,\dots,T\)，\(\theta_t\) 单射。
\end{enumerate}
\end{theorem}
\begin{proof}
对有限维滤过空间之间的严格滤过线性映射，若其关联分级映射单射，则原映射单射；反之亦然。此为滤过线性代数基本定理，因各 \(K_t\) 有限维。
\end{proof}

此定理将观测问题转化为对每一层首次观测映射 \(\theta_t\) 的单射性分析。

\section{形式首次观测与 leading term 理论}

\subsection{极值路径与形式首次观测}
对每个基索引 \(i\in A_T\)，若 \(\tau(i)\le T\)，则其首次接触窗口的极值路径唯一确定一个“最快”非零观测候选。

\begin{definition}[极值接触位置]
设 \(\tau(i)\le T\)。
\begin{itemize}
    \item 若 \(i<-cT\)（左侧进入），则 \(d_+>0\) 且 \(\tau(i)=\lceil(-cT-i)/d_+\rceil\)。定义极值接触位置 \(j_*(i) = -cT\)，对应单项式系数 \(c_*(i) = a_{d_+}^{\tau(i)} \neq 0\)。
    \item 若 \(i>cT\)（右侧进入），则 \(d_->0\) 且 \(\tau(i)=\lceil(i-cT)/d_-\rceil\)。定义极值接触位置 \(j_*(i) = cT\)，对应单项式系数 \(c_*(i) = a_{-d_-}^{\tau(i)} \neq 0\)。
    \item 若 \(|i|\le cT\)，则 \(\tau(i)=0\)，定义 \(j_*(i)=i\)，\(c_*(i)=1\)。
\end{itemize}
\end{definition}

极值路径的唯一性来自极值位移的最大性：任何其他路径在同一时间内净位移严格小于极值，无法在该时间触及窗口边界。因此，在 \(\pi F^{\tau(i)} e_i\) 中，位于 \(j_*(i)\) 的系数恰为 \(c_*(i)\)，且该位置不可能由其他更短或等长的路径以同等净位移到达（除非另有相等极值位移的单步系数，此时会出现多个等速路径，但这是 \(r\ge2\) 且多系数非零的情形，稍后讨论）。不失一般性，我们称 \((t=\tau(i), j=j_*(i))\) 为 \(e_i\) 的\emph{leading term 位置}，其系数为 \(c_*(i)\)。

\begin{definition}[形式首次观测映射]
对每个 \(i\) 满足 \(\tau(i)=t\le T\)，定义向量
\[
\ell_i = \theta_t(e_i + K_t) = \pi F^t e_i \in \F_q^{[-cT,cT]}.
\]
该向量在位置 \(j_*(i)\) 上的分量为 \(c_*(i)\)。若 \(c_*(i)\neq0\)，则 \(\ell_i\) 非零，且 \(e_i\notin K_t\)，即 \(\nu(e_i)=t\)。
\end{definition}

若对某 \(i\) 有 \(c_*(i)=0\)，则由于有限域特征可能发生极端稀疏化，此时 \(e_i\in K_{\tau(i)}\)，首次非零观测时间推迟。本文主要关注 \(c_*(i)\neq0\) 的情形；Frobenius 诱导的零化将单独讨论。

\subsection{Leading term 矩阵与分级单射性}
固定时间层 \(t\)，设 \(J_t = \{i\in I_t : c_*(i)\neq0\}\)。将 \(E_t\) 的生成元（不一定是基）\(\{e_i + K_t\}_{i\in J_t}\) 按 \(j_*(i)\) 的某种顺序排列，并考察 \(\theta_t\) 在该生成元族上的矩阵。

\begin{lemma}[下三角结构]\label{lem:theta_t_triangular}
设 \(i,i'\in J_t\) 且 \(\tau(i)=\tau(i')=t\)。
\begin{enumerate}[label=(\roman*)]
    \item 若 \(j_*(i) = j_*(i')\)，则 \(\ell_i\) 与 \(\ell_{i'}\) 在位置 \(j_*(i)\) 上的系数分别为 \(c_*(i),c_*(i')\)。
    \item 若 \(j_*(i) \neq j_*(i')\)，且按左右边界分为不同侧，则 \(\Supp(\ell_i)\) 与 \(\Supp(\ell_{i'})\) 在对方极值位置上的交集为零。
    \item 若 \(j_*(i),j_*(i')\) 位于同一侧边界（例如均为右边界 \(cT\)），且 \(i < i'\)，则 \(j_*(i')\) 在 \(\ell_i\) 中的系数为零。
\end{enumerate}
\end{lemma}
\begin{proof}
由极值路径的位移唯一性与窗口外区域无观测，直接验证即得。详见附录~\ref{sec:leading_term_proof}。
\end{proof}

该引理表明，若按 \(j_*(i)\) 的自然序排列生成元，\(\theta_t\) 在该族上的矩阵呈分块下三角形态，对角元为 \(c_*(i)\)。当所有 \(c_*(i)\) 非零且 leading term 位置两两不同时，该矩阵可逆，从而 \(\theta_t\) 单射。

\begin{theorem}[Leading term 单射准则]\label{thm:leading_term_injective}
若对每个 \(t=0,\dots,T\)，集合 \(\{ (\tau(i), j_*(i)) : i\in A_T,\ \tau(i)=t,\ c_*(i)\neq0\}\) 中的元素两两不同，且每个 \(c_*(i)\neq0\)，则 \(\gr(\Theta_T)\) 单射，从而 \(H_T = \dim G_T\)（即代数核为零，观测缺陷完全由几何核贡献）。
\end{theorem}
\begin{proof}
在上述条件下，每个 \(\theta_t\) 在生成元族上的矩阵为可逆下三角矩阵，故 \(\theta_t\) 是单射。由定理~\ref{thm:associated_graded}，\(\gr(\Theta_T)\) 单射，进而 \(\Theta_T\) 在 \(V_0/G_T\) 上单射，代数核为零。
\end{proof}

该定理将极值满秩子矩阵（命题~\ref{prop:extreme_submatrix}）作为特例纳入：极值满秩子矩阵的列恰选自那些 leading term 互异的极值索引，它们保证了对应分块的可逆性。但定理~\ref{thm:leading_term_injective} 更进一步，指出了整个分级系统单射的充分条件。

\subsection{碰撞图与代数核的表述}
当定理条件不满足时，多个基向量可能共享同一 leading term 位置 \((t,j)\)，或虽位置不同但因多极值项同时存在而产生更复杂的线性关系。为此引入碰撞图。

\begin{definition}[Leading term 碰撞图]
定义有向二部图 \(\mathcal{G}\)：顶点集为 \(A_T\)（初始索引）与观测时空点 \(\{(t,j): 0\le t\le T, |j|\le cT\}\)。当 \(\tau(i)=t\) 且 \(c_*(i)\neq0\) 时，连边 \(i \to (t, j_*(i))\)，权重为 \(c_*(i)\)。
若两个不同的 \(i,i'\) 指向同一 \((t,j)\)，则称发生\emph{一级碰撞}。
\end{definition}

一级碰撞意味着在对应的 \(\theta_t\) 矩阵中，两列的首个非零元处于同一行且系数已知。此时可能通过线性组合消去 leading term，产生属于 \(\ker\theta_t\) 的非零元素，从而贡献代数核。更复杂的多级碰撞（通过中间层传播）原则上都可归结为碰撞图的组合结构。

\begin{proposition}[碰撞诱导的缺陷界]\label{prop:collision_bound}
设一级碰撞次数为 \(N_{\mathrm{coll}}\)，即
\[
N_{\mathrm{coll}} = \sum_{t=0}^T \bigl( |J_t| - |\{j_*(i): i\in J_t\}| \bigr).
\]
则
\[
\dim\ker(\gr(\Theta_T)) \ge N_{\mathrm{coll}},
\]
因而 \(H_T \le \dim G_T + (n_T - \dim G_T) - N_{\mathrm{coll}}\)（上界）且 \(\dim K_T^{\mathrm{alg}} \ge N_{\mathrm{coll}}\)（下界，若未进一步抵消）。
\end{proposition}
\begin{proof}
在 \(\theta_t\) 的矩阵中，每出现一对 \(i\neq i'\) 指向同一 \(j\)，列向量在 leading 位置上的系数向量成比例（比例 \(c_*(i)/c_*(i')\)），由此可构造一个线性相关的向量对。这些关系线性独立时即贡献核维数。
\end{proof}

\textbf{注：} 一级碰撞不一定导致真正线性相关（若后续行可区分），但结合下三角结构，当仅考虑 leading 行时它们确实线性相关。因此代数核的存在性紧密联系于碰撞图中无法通过更深层观测消除的依赖关系。

\section{极值满秩子矩阵与 \(r=1\) 完全可观测}
\begin{proposition}[极值满秩子矩阵——局部 pivot 证书]\label{prop:extreme_submatrix}
设 \(a_{-d_-}\neq0\) 且 \(a_{d_+}\neq0\)。则观测矩阵 \(\Theta_T\) 包含一个 \((2cT+1 + N_- + N_+)\times(2cT+1 + N_- + N_+)\) 的可逆子矩阵，其中 \(N_\pm\) 为有效极值层数（定义见附录~\ref{sec:card_I}）。因而
\[
R_T \ge 2cT+1 + N_- + N_+,\qquad
H_T \le 2rT - N_- - N_+.
\]
\end{proposition}
\begin{proof}
该子矩阵的列恰为 leading term 互异的极值索引集，其 leading term 矩阵可逆。详细构造同原文，不赘述。
\end{proof}

当 \(r=1\) 时，\(d_-=d_+\in\{0,1\}\)。若两端系数非零，则 \(d_-=d_+=1=r\)，几何核消失，且 \(N_-=N_+=\min(T,T)=T\)。此时所有基向量的 leading term 两两不同：中心部分 \(\tau=0\) 各列 \(j_*(i)=i\) 均不同；左侧极值列的 \(j_*(i)=-cT\) 且对应时间 \(\tau\) 递增，故 \((t,j)\) 唯一；右侧类似。因此定理~\ref{thm:leading_term_injective} 的条件完全满足，立得完全可观测性。

\begin{theorem}[\(r=1\) 完全可观测性]\label{thm:r1_full_obs}
设 \(r=1\)，\(a_{-1}\neq0\) 且 \(a_1\neq0\)。则对任意 \(T\ge0\) 和 \(c\ge1\)，\(H_T=0\)。
\end{theorem}
\begin{proof}
此时 leading term 全部互异，由定理~\ref{thm:leading_term_injective} 得 \(\gr(\Theta_T)\) 单，从而 \(\Theta_T\) 列满秩。
\end{proof}

\section{代数核与 Frobenius 风险}
当 \(r\ge2\) 且 \(a_{-r},a_r\neq0\) 时，几何核消失，\(H_T\le 2T(r-1)\)。代数核的可能来源有三类：
\begin{enumerate}[label=(\arabic*)]
    \item \textbf{Frobenius 稀疏化：} 在特征 \(p\) 的域上，\(P(z)^{p^m}\) 可能发生二项式展开的稀疏化，致使某些极值系数为零，\(c_*(i)=0\)，从而 leading term 直接消失，对应的 \(e_i\) 首次非零观测延迟。
    \item \textbf{Leading term 碰撞：} 即使系数非零，不同的 \(i,i'\) 可能共享同一 \((t,j_*(i))\)，产生一级碰撞。后续观测能否将其分离，取决于高阶项的系数结构。
    \item \textbf{深层 Toeplitz syzygy：} 即使 leading term 完全分离，多项式矩方程可能因卷积码式的内部关系产生非平凡核。此为最隐蔽的代数障碍。
\end{enumerate}

\begin{openproblem}[几何版——无 Frobenius 共振时的零缺陷]\label{prob:full_obs_geo}
设 \(c\ge r\)，\(a_{-r}\neq0\) 且 \(a_r\neq0\)。若规则多项式 \(P(z)\) 的幂次在时间 \(\le T\) 内未因域特征发生非平凡稀疏化（即对所有 \(t\le T\)，\(c_*(i)\neq0\) 恒成立），且 leading term 碰撞图为一对一匹配，是否必有 \(H_T=0\)？
\end{openproblem}

\begin{openproblem}[有限域版——Frobenius 共振诱导的反例]\label{prob:full_obs_frob}
在 \(\F_{p^k}\) 上，是否存在参数 \(a_j\) 及 \(T\)，使得 \(a_{-r},a_r\neq0\) 但 \(H_T>0\)？特别地，特征诱导的 \(c_*(i)=0\) 或碰撞后不可分离，是否生成非平凡有限支撑 syzygy？
\end{openproblem}

\begin{conjecture}[有限域反例存在性]\label{conj:frob_counter}
在 \(\F_2\) 上，存在 \(r\ge2\) 的规则（例如 \(P(z)=z^{-2}+1+z^2\)）及足够大的 \(T\ge 2\)，使得 Frobenius 稀疏化导致某些极值系数为零并产生跨时间抵消，从而 \(H_T>0\)。
\end{conjecture}

\textbf{注：} 该猜想目前仅有低时间深度的数值迹象，未经一般性证明。碰撞图与 leading term 矩阵的分析是判断这类反例是否存在的最直接工具。

\section{与 Toeplitz 核及卷积码的关系}
观测缺陷等价于截断 Toeplitz 算子束的核问题。记 \(\Pi_{cT}: \F_q[z,z^{-1}]\to \F_q^{[-cT,cT]}\) 为取系数映射，则
\[
\ker\Theta_T \cong \bigl\{ X\in\F_q[z,z^{-1}] : \Supp(X)\subseteq A_T,\;
\Pi_{cT}(P^t X)=0 \ \forall\, t\le T \bigr\}.
\]
leading term 分析对应 Laurent 多项式乘法的边界行为，碰撞图反映多项式系数的交互。卷积码中类似结构出现于有限约束长度序列的代数性质，本框架为其提供了显式可计算的几何–代数双层分析工具。

\section{结论}
本文建立了几何过滤与关联分级的统一框架，将线性元胞自动机的有限时空可观测性归结为分级映射 \(\gr(\Theta_T)\) 的单射性。通过定义形式首次观测的 leading term，证明当 leading term 两两不同时系统完全可观（定理~\ref{thm:leading_term_injective}）。该结论严格包含 \(r=1\) 非退化情形的完全可观测性，并给出了一般情形下代数核的组合刻画（碰撞图）。几何核维数由极值步幅显式给出，所有参数均可计算。尚未解决的代数核存在性问题已明确拆分为无 Frobenius 共振的几何版与有限域 Frobenius 共振版，并指出 leading term 碰撞分析是解决这些问题的核心路径。

\appendix
\section{关键证明细节}
\subsection{\(\tau(i)\) 公式与 \(I_\ell\) 基数}\label{sec:card_I}
\(\tau(i)\) 的显式公式由定义直接导出。对于右半侧：
\[
I_\ell = \{ i : \lceil(i-cT)/d_-\rceil = \ell \} \cap A_T
= \{ i : cT+(\ell-1)d_- < i \le cT+\ell d_-,\; |i|\le (c+r)T \}.
\]
故 \(|I_\ell| = \min(cT+\ell d_-,\,(c+r)T) - \max(cT+(\ell-1)d_-+1,\,cT+1) +1\)，通常为 \(d_-\)，仅在触及 \(A_T\) 边界时截断。极值列 \(i=cT+d_-\ell\) 存在于 \(A_T\) 当且仅当 \(\ell\le \lfloor rT/d_-\rfloor\) 且 \(\ell\le T\)，故总数为 \(N_-=\min(T,\lfloor rT/d_-\rfloor)\)。左半侧对偶定义 \(N_+=\min(T,\lfloor rT/d_+\rfloor)\)。

\subsection{引理~\ref{lem:theta_t_triangular} 的证明}\label{sec:leading_term_proof}
固定 \(t\)，设 \(i,i'\in J_t\)。
\begin{enumerate}[label=(\roman*)]
    \item 若 \(j_*(i)=j_*(i')\)，按定义两者 leading term 处于同一观测位置，系数各为 \(c_*(i),c_*(i')\)。
    \item 若分属左右不同边界，例如 \(j_*(i)=cT\)，\(j_*(i')=-cT\)，则 \(\Supp(\ell_i)\subseteq (-\infty, cT]\) 且最右端为 \(cT\)；\(\Supp(\ell_{i'})\subseteq [-cT,\infty)\) 且最左端为 \(-cT\)。因 \(cT>0\)，交集在对方极值位置上的系数为零。
    \item 若同侧，例如均为右侧 \(cT\)，且 \(i<i'\)。时刻 \(t\) 的观测由 \(F^t e_i\) 给出，其支撑右端点为 \(i+d_+ t\)。当 \(i' = i + k d_-\)（满足进入条件）时，\(j_*(i')=cT\)。由于 \(i'>i\)，\(F^t e_i\) 的支撑右端点 \(i+d_+ t < i'\)，不可能到达 \(cT\) 位置（因为要到达 \(cT\) 需要净位移更大），故在 \(j_*(i')\) 处系数为零。这就保证了下三角性。
\end{enumerate}
更正式的论证使用极值路径的唯一性和非极值路径位移严格小于极值位移这一事实，详见规则多项式的极值分解。

\subsection{规则90的边界单调漂移}\label{sec:rule90}
对 \(\F_2\) 上规则90，\(P(z)=z+z^{-1}\)，\(d_-=d_+=1\)。由 \(\tau\) 公式得 \(V_\ell=\spanop\{e_{-cT-\ell},\,e_{cT+\ell}\}\) (\(\ell\ge1\))。
计算
\[
F^\ell e_{-cT-\ell} = \sum_{k=0}^{\ell} \binom{\ell}{k} e_{-cT-2k}\pmod 2.
\]
当 \(k=0\) 得 \(e_{-cT}\)；\(k\ge1\) 时指数 \(\le -cT-2<-cT\)，完全落在窗口外。故 \(\pi F^\ell e_{-cT-\ell} = e_{-cT}\)。
同理 \(\pi F^\ell e_{cT+\ell}=e_{cT}\)。因此在分级块上，两个基向量分别映为标准基向量，leading term 互异，定理~\ref{thm:leading_term_injective} 的条件成立，观测矩阵完全解耦为若干低维可逆块的直和。

\section{代数核消失定理：几何核的完备性}

我们现给出一个严谨的归纳证明，确认在极值方向均非零的条件下，核过滤与几何过滤完全重合，从而观测缺陷仅由几何核贡献，代数核恒为零。

\subsection{设定与已知事实}
沿用前文的全部记号：规则多项式 \(P(z)=\sum_{j=-r}^r a_jz^j\)，\(d_-=\max\{|j|:j<0,\,a_j\neq 0\}\)，\(d_+=\max\{j:j>0,\,a_j\neq 0\}\)，若某侧无系数则对应值取 \(0\)。观测窗口 \(\Lambda_T=\{(t,i):0\le t\le T,\,|i|\le cT\}\)，\(c\ge r\)。初始空间 \(V_0=\spanop\{e_i:|i|\le (c+r)T\}\)，几何接触时间 \(\tau(i)\) 如定义，几何过滤 \(G_t=\spanop\{e_i:\tau(i)>t\}\)，核过滤 \(K_t=\bigcap_{s=0}^t\ker(\pi F^s)\)。已知 \(G_t\subseteq K_t\) 对所有 \(t\) 成立。我们假设 \(d_->0\) 且 \(d_+>0\)（两端至少各有一个非零系数），且对应的极值系数 \(a_{-d_-},a_{d_+}\) 非零。

\subsection{基情形：\(t=0\)}
由于 \(\tau(i)=0\) 当且仅当 \(|i|\le cT\)，故 \(G_0=\spanop\{e_i:|i|>cT\}\)。而 \(K_0=\ker\pi\) 恰为在坐标窗口 \([-cT,cT]\) 上恒零的向量，自然等于 \(G_0\)。因此
\begin{equation}
  K_0=G_0.
\end{equation}

\subsection{归纳步骤：从 \(K_{t-1}=G_{t-1}\) 到 \(K_t=G_t\)}
假设已对某个 \(t\ge 1\) 证明了 \(K_{t-1}=G_{t-1}=\spanop\{e_i:\tau(i)\ge t\}\)（这里 \(\tau(i)\ge t\) 包含 \(\infty\)）。考虑映射
\[
  L_t: K_{t-1} \to \F_q^{[-cT,cT]},\quad L_t(v)=\pi F^t v.
\]
则 \(K_t = K_{t-1}\cap\ker(\pi F^t)=\ker L_t\)。
由于 \(K_{t-1}\) 由基 \(\{e_i:\tau(i)\ge t\}\) 张成，且当 \(\tau(i)>t\) 时显然 \(e_i\in G_t\subseteq K_t\subseteq\ker L_t\)，因此 \(G_t\subseteq\ker L_t\)。我们只需证明 \(\ker L_t = G_t\)。将 \(K_{t-1}\) 中的向量按 \(\tau\) 值分解：
\[
  v = \sum_{i\in I_t} c_i e_i + \sum_{j:\tau(j)>t} c_j e_j,
\]
其中 \(I_t=\{i\in A_T:\tau(i)=t\}\)。则 \(L_t(v)=\sum_{i\in I_t} c_i \pi F^t e_i\)，因为对 \(\tau(j)>t\) 有 \(\pi F^t e_j=0\)（几何包含）。故 \(v\in\ker L_t\) 当且仅当 \(\sum_{i\in I_t} c_i u_i =0\)，其中 \(u_i:=\pi F^t e_i\)。
若能证明向量组 \(\{u_i\}_{i\in I_t}\) 线性无关，则 \(v\in\ker L_t\) 推出所有 \(c_i=0\)，从而 \(v\) 仅由 \(\tau(j)>t\) 的基生成，即 \(v\in G_t\)。于是 \(\ker L_t=G_t\)，归纳完成。

因此，整个问题归结为证明下述命题。

\begin{lemma}[同层观测向量的线性无关性]\label{lem:indep}
设 \(d_->0\) 且 \(d_+>0\)。对任意 \(t\ge 1\)，向量组
\[
  \{u_i = \pi F^t e_i : i\in I_t\}
\]
在 \(\F_q^{[-cT,cT]}\) 中线性无关。
\end{lemma}

\begin{proof}
将 \(I_t\) 划分为右侧进入和左侧进入两部分：
\begin{align*}
  I_t^R &= \{i\in I_t : i > cT\},\\
  I_t^L &= \{i\in I_t : i < -cT\}.
\end{align*}
（注意 \(\tau(i)=t\ge 1\) 时必有 \(|i|>cT\)，故 \(I_t=I_t^R\cup I_t^L\)。）

对 \(i\in I_t^R\)，由极值位移定义，其支撑满足
\[
  \Supp(\pi F^t e_i) \subseteq [j_*(i),\, \min(cT,\, i+d_+t)],
\]
其中 \(j_*(i)=i-d_-t\)。因为 \(\tau(i)=t\) 是最小首次进入时间，易见 \(cT-d_- < j_*(i)\le cT\)，故 \(j_*(i)\ge cT-d_- \ge 0\)（因 \(c\ge r\ge d_-\)）。从而所有 \(u_i\) (\(i\in I_t^R\)) 的支撑完全包含于非负半轴 \([0,cT]\)。
同理，对 \(i\in I_t^L\)，令 \(j_*(i)=i+d_+t\)，其满足 \(-cT\le j_*(i)\le -cT+d_+ \le 0\)，且支撑完全含于非正半轴 \([-cT,0]\)。

因此，右侧向量的支撑与左侧向量的支撑仅在可能包含的 \(0\) 处有交集。现按如下方式证明线性无关。

首先，将观测行划分为非负行 \(R_+=\{j:0\le j\le cT\}\) 和负行 \(R_-=\{-cT,\dots,-1\}\)。在负行上，所有 \(u_i\ (i\in I_t^R)\) 均为零；在正行上，所有 \(u_i\ (i\in I_t^L)\) 均为零（0 行归入某一边，不妨归入 \(R_+\)）。构造列向量组的表示矩阵，行按先 \(R_-\) 后 \(R_+\) 排列，列按先 \(I_t^L\) 后 \(I_t^R\) 排列。此时矩阵呈分块对角形：
\[
  M = \begin{pmatrix}
    M_L & 0 \\[2pt]
    0 & M_R
  \end{pmatrix},
\]
其中 \(M_L\) 的行取自 \(R_-\)（列 \(I_t^L\)），\(M_R\) 的行取自 \(R_+\)（列 \(I_t^R\)）。因为左侧列在正行全零，右侧列在负行全零。因此，\(\{u_i\}\) 线性无关当且仅当 \(M_L\) 的列线性无关且 \(M_R\) 的列线性无关。

现分别处理右侧和左侧。
\paragraph{右侧子矩阵 \(M_R\) 的可逆性。}
行取 \(R_+\) 中所有满足 \(j=j_*(i)\) 的行，其中 \(i\in I_t^R\)。因 \(j_*(i)\) 互异（若 \(i<i'\)，则 \(i-d_-t < i'-d_-t\)），这些行构成一个 \(|I_t^R|\) 维子集。将 \(I_t^R\) 中的指标按 \(j_*(i)\) 递增的顺序排列，相应的行也按相同顺序排列，得到一个 \(|I_t^R|\times |I_t^R|\) 子矩阵 \(\tilde M_R\)。我们断言 \(\tilde M_R\) 是下三角矩阵且对角元非零。

取 \(i,i'\in I_t^R\) 且 \(j_*(i)<j_*(i')\)（即 \(i<i'\)）。在行 \(j_*(i)\) 处，列 \(i'\) 的元素是 \(\pi F^t e_{i'}\) 在 \(j_*(i)\) 的系数。由于 \(i'\) 在 \(t\) 步内的支撑左边界为 \(j_*(i')\)，而 \(j_*(i)<j_*(i')\)，因此该位置系数为零。从而 \(\tilde M_R\) 在 \(j_*(i)<j_*(i')\) 时为零，即严格下三角部分为零。对角元为 \(u_i\) 在 \(j_*(i)\) 处的值。由极值路径的唯一性，该值等于 \(a_{-d_-}^t\)，非零。故 \(\tilde M_R\) 是可逆的下三角矩阵，从而 \(M_R\) 的列（在原始行集合下）线性无关。

\paragraph{左侧子矩阵 \(M_L\) 的可逆性。}
对称地，取行 \(R_-\) 中所有 \(j=j_*(i)\) (\(i\in I_t^L\))，按 \(j_*(i)\) 递减的顺序排列（因左侧极值位置越右越大，按绝对值递减）。类似论证给出一个上三角（或适当排序下的下三角）可逆矩阵，对角元为 \(a_{d_+}^t\neq 0\)。细节从略。

由于 \(M_L\) 和 \(M_R\) 均列满秩，分块对角矩阵 \(M\) 的列线性无关，故 \(\{u_i\}\) 线性无关。
\end{proof}

\subsection{核过滤的完全确定}
由引理~\ref{lem:indep}，我们立即得到对每个 \(t\ge 1\)，
\[
  K_t = \ker L_t = \spanop\{e_i:\tau(i)>t\} = G_t.
\]
由归纳法，此式对所有 \(t=0,1,\dots,T\) 成立。特别地，
\[
  \ker\Theta_T = K_T = G_T.
\]
因此我们证明了以下定理。

\begin{theorem}[代数核消失与核的几何完全性]\label{thm:final}
设线性元胞自动机的规则满足 \(d_->0\) 且 \(d_+>0\)（即左、右两个方向均存在非零系数）。则对任意观测时长 \(T\ge 0\) 和窗口参数 \(c\ge r\)，核过滤与几何过滤完全重合：
\[
  K_t = G_t,\qquad t=0,1,\dots,T.
\]
等价地，\(\ker\Theta_T = G_T\)，观测缺陷完全由几何核贡献：
\[
  H_T = \dim G_T = T(r-d_-) + T(r-d_+).
\]
特别地，若极端系数均非零（\(a_{-r}\neq0,\ a_r\neq0\)），则 \(d_-=d_+=r\)，\(H_T=0\)，即系统完全可观测。
\end{theorem}

\subsection{对先前开放问题与猜想的回应}
定理~\ref{thm:final} 彻底解决了前文提出的全部核心问题：
\begin{itemize}
  \item \textbf{Leading term 完备性}：不仅没有碰撞产生新核，实际上连碰撞的可能性都不存在，因为同层不同基向量的极值观测位置自动互异且左右分离，构成可逆分块对角阵。所有核均来自几何上完全不可达的模态。
  \item \textbf{Frobenius 稀疏化}：有限域的特征不会使极值系数 \(a_{\pm d_\pm}^t\) 变为零，因此极值路径的非零性与唯一性始终保留，代数核不可能因 Frobenius 而产生。
  \item \textbf{Toeplitz syzygy 与隐藏核}：不存在所谓的“隐形 syzygy”。核过滤在每一步均与几何过滤相等，无任何额外的代数抵消空间。
\end{itemize}

至此，因果窗口下线性元胞自动机的可观测性问题获得完全分类：缺陷仅由最大传播速度的不足造成，形式由极值位移 \(d_\pm\) 显式给出。结构分解 \(\ker\Theta_T = G_T\) 是严格的，且证明不依赖任何未验证假设。

\end{document}